\begin{lemma}
\label{lemma-additivity-loc-chern-c}
Let $(S, \delta)$ be as in Situation \ref{situation-setup}.
Let $X$ be locally of finite type over $S$. Let $Z \to X$ be
a closed immersion. Let
$$
E_1 \to E_2 \to E_3 \to E_1[1]
$$
be a distinguished triangle of perfect objects in $D(\mathcal{O}_X)$.
Assume
\begin{enumerate}
\item the restrictions $E_1|_{X \setminus Z}$ and $E_3|_{X \setminus Z}$
are isomorphic to finite locally free $\mathcal{O}_{X \setminus Z}$-modules
of rank $< p_1$ and $< p_3$ placed in degree $0$, and
\item at least one of the following is true:
(a) $X$ is quasi-compact,
(b) $X$ has quasi-compact irreducible components,
(c) $E_3 \to E_1[1]$ can be represented by a map of locally
bounded complexes of finite locally free $\mathcal{O}_X$-modules, or
(d) there exists an envelope $f : Y \to X$ such that $Lf^*E_3 \to Lf^*E_1[1]$
can be represented by a map of locally bounded complexes of
finite locally free $\mathcal{O}_Y$-modules.
\end{enumerate}
With notation as in Remark \ref{remark-loc-chern-classes} we have
$$
c^{(p_1 + p_3)}(Z \to X, E_2) = c^{(p_1)}(Z \to X, E_1)c^{(p_3)}(Z \to X, E_3)
$$
in $A^{(p_1 + p_3)}(Z \to X)$.
\end{lemma}

\begin{proof}
Observe that the assumptions imply that $E_2|_{X \setminus Z}$ is zero,
resp.\ isomorphic to a finite locally free $\mathcal{O}_{X \setminus Z}$-module
of rank $< p_1 + p_3$. Thus the statement makes sense.

\medskip\noindent
Let $f : Y \to X$ be an envelope. Expanding the left and right hand sides
of the formula in the statement of the lemma we see that we have to prove
some equalities of classes in $A^*(X)$ and in $A^*(Z \to X)$. By the
uniqueness in Lemma \ref{lemma-envelope-bivariant} it suffices to prove the
corresponding relations in $A^*(Y)$ and $A^*(Z \to Y)$. Since moreover
the construction of the classes involved is compatible with base change
(Lemma \ref{lemma-base-change-loc-chern}) we may replace $X$ by $Y$
and the distinguished triangle by its pullback.

\medskip\noindent
In the proof of Lemma \ref{lemma-additivity-on-perfect} we have
seen that conditions (2)(a), (2)(b), and (2)(c) imply condition
(2)(d). Combined with the discussion in the previous paragraph we
reduce to the case discussed in the next paragraph.

\medskip\noindent
Let $\varphi^\bullet : \mathcal{E}_3^\bullet[-1] \to \mathcal{E}_1^\bullet$
be a map of locally bounded complexes of finite locally free
$\mathcal{O}_X$-modules representing the map $E_3[-1] \to E_1$
in the derived category. Consider the scheme
$X' = \mathbf{A}^1 \times X$ with projection
$g : X' \to X$. Let $Z' = g^{-1}(Z) = \mathbf{A}^1 \times Z$.
Denote $t$ the coordinate on $\mathbf{A}^1$. Consider the cone
$\mathcal{C}^\bullet$ of the map of complexes
$$
t g^*\varphi^\bullet :
g^*\mathcal{E}_3^\bullet[-1]
\longrightarrow
g^*\mathcal{E}_1^\bullet
$$
over $X'$. We obtain a distinguished triangle
$$
g^*\mathcal{E}_1^\bullet \to \mathcal{C}^\bullet \to
g^*\mathcal{E}_3^\bullet \to g^*\mathcal{E}_1^\bullet[1]
$$
where the first three terms form a termwise split short exact
sequence of complexes. Clearly $\mathcal{C}^\bullet$ is a
bounded complex of finite locally free $\mathcal{O}_{X'}$-modules
whose restriction to $X' \setminus Z'$ is isomorphic to a
finite locally free
$\mathcal{O}_{X' \setminus Z'}$-module of rank $< p_1 + p_3$
placed in degree $0$. Thus we have the localized Chern classes
$$
c_p(Z' \to X', \mathcal{C}^\bullet) \in A^p(Z' \to X')
$$
for $p \geq p_1 + p_3$. For any $\alpha \in \CH_k(X)$ consider
$$
c_p(Z' \to X', \mathcal{C}^\bullet) \cap g^*\alpha
\in \CH_{k + 1 - p}(\mathbf{A}^1 \times X)
$$
If we restrict to $t = 0$, then the map $t g^*\varphi^\bullet$
restricts to zero and $\mathcal{C}^\bullet|_{t = 0}$
is the direct sum of $\mathcal{E}_1^\bullet$ and $\mathcal{E}_3^\bullet$.
By compatibility of localized Chern classes with base change
(Lemma \ref{lemma-base-change-loc-chern}) we conclude that
$$
i_0^* \circ c^{(p_1 + p_3)}(Z' \to X', \mathcal{C}^\bullet) \circ g^* =
c^{(p_1 + p_2)}(Z \to X, E_1 \oplus E_3)
$$
in $A^{(p_1 + p_3)}(Z \to X)$. On the other hand, if we restrict to $t = 1$,
then the map $t g^*\varphi^\bullet$
restricts to $\varphi$ and $\mathcal{C}^\bullet|_{t = 1}$
is a bounded complex of finite locally free modules representing $E_2$.
We conclude that
$$
i_1^* \circ c^{(p_1 + p_3)}(Z' \to X', \mathcal{C}^\bullet) \circ g^* =
c^{(p_1 + p_2)}(Z \to X, E_2)
$$
in $A^{(p_1 + p_3)}(Z \to X)$. Since $i_0^* = i_1^*$ by definition of
rational equivalence (more precisely this follows from the formulae in
Lemma \ref{lemma-linebundle-formulae}) we conclude that
$$
c^{(p_1 + p_2)}(Z \to X, E_2) = c^{(p_1 + p_2)}(Z \to X, E_1 \oplus E_3)
$$
This reduces us to the case discussed in the next paragraph.

\medskip\noindent
Assume $E_2 = E_1 \oplus E_3$ and the triples $(X, Z, E_i)$
are as in Situation \ref{situation-loc-chern}.
For $i = 1, 3$ let
$$
b_i : W_i \to \mathbf{P}^1_X
\quad\text{and}\quad
Q_i
\quad\text{and}\quad
T'_i \subset T_i \subset W_{i, \infty}
$$
be as in the proof of Lemma \ref{lemma-independent-loc-chern}.
By definition
$$
c_p(Z \to X, E_i) = c'_p(Q_i)
$$
where the right hand side is the bivariant class constructed in
Lemma \ref{lemma-localized-chern-pre} using $W_i, b_i, Q_i, T'_i$.
Set $W = W_1 \times_{b_1, \mathbf{P}^1_X, b_2} W_2$ and consider
the cartesian diagram
$$
\xymatrix{
W \ar[d]_{g_1} \ar[rd]^b \ar[r]_{g_3} & W_3 \ar[d]^{b_3} \\
W_1 \ar[r]^{b_1} & \mathbf{P}^1_X
}
$$
Of course $b^{-1}(\mathbf{A}^1)$ maps isomorphically to $\mathbf{A}^1_X$.
Observe that $T' = g_1^{-1}(T'_1) \cap g_2^{-1}(T'_2)$ still contains
all the points of $W_\infty$ lying over $X \setminus Z$.
By Lemma \ref{lemma-localized-chern-pre-independent} we may use
$W$, $b$, $g_i^*\mathcal{Q}_i$, and
$T'$ to construct $c_p(Z \to X, E_i)$ for $i = 1, 3$.
Also, by the stronger independence given in
Lemma \ref{lemma-independent-loc-chern-bQ} we may use
$W$, $b$, $g_1^*Q_1 \oplus g_3^*Q_3$, and $T'$
to compute the classes $c_p(Z \to X, E_2)$.
Thus the desired equality follows from
Lemma \ref{lemma-localized-chern-pre-sum-c}.
\end{proof}